\section{Tổng kết và các hướng nghiên cứu}

\subsection{Tổng kết}

\begin{itemize}
  \item \textbf{PCA} tìm chiếu tuyến tính bậc $k$ nhỏ nhất lỗi tái tạo Frobenius; nghiệm gắn với vector kỳ dị của ma trận hiệp phương sai $\mathbf{C}$ (Định lý~\ref{thm:pca}).
  \item \textbf{KPCA} làm việc với ma trận Gram $\mathbf{K}$, cho phép phi tuyến mà không cần $\Phi$ tường minh.
  \item \textbf{Isomap, Laplacian eigenmaps, LLE} xây kernel hoặc ma trận Laplacian từ đồ thị láng giềng; bảng~\ref{tab:dr-compare} so sánh mức ``toàn cục'' hay ``cục bộ'' mà mỗi phương pháp ưu tiên giữ.
  \item \textbf{Johnson--Lindenstrauss} cho thấy chiều $k = O(\log m / \varepsilon^2)$ đủ để bảo toàn khoảng cách cặp dưới chiếu ngẫu nhiên.
  \item \textbf{[MỞ RỘNG] t-SNE} và \textbf{UMAP} (Mục~\ref{sec:extension-viz}): tối ưu nhúng phi tuyến qua KL hoặc cross-entropy fuzzy trên đồ thị láng giềng; phù hợp trực quan hóa hơn eigenmaps tuyến tính thuần túy trong sách.
\end{itemize}

\paragraph{Hạn chế cần nhớ khi áp dụng.}
Manifold lý tưởng (kiểu Swiss roll) hiếm khớp dữ liệu thật; $\mathbf{K}_{\mathrm{Iso}}$ có thể không xác định dương bán trên mẫu hữu hạn; Laplacian và LLE nhạy với số láng giềng $t$, tham số $\sigma$ và cách nối cạnh; PCA/KPCA tuyến tính (kernel đơn giản) không tự khắc phục cấu trúc cong phức tạp nếu không chọn kernel hoặc thuật toán phù hợp.

\subsection{Các hướng nghiên cứu}

Khảo sát chi tiết các bài báo sau năm 2021 nằm ở Chương~\ref{sec:research-contrastive} (Một số hướng nghiên cứu). Các paper sẽ được liên hệ với lý thuyết sách (PCA, manifold, JL) và với UMAP \textbf{[MỞ RỘNG]} (Mục~\ref{sec:umap}).
